The complete charged energy measurement also controls position for its
specified preparation.  Write $H=H_0+\lambda T+4N_{\rm ref}$ for the
eight-cell interacting source, with $\lambda=2^{-34}$ and
$\|\lambda T\|\leq v=360\,2^{-34}$.  The reference is a retained static
charge at the origin.  Let $\chi=P_I(H)B_r\Psi$ be the unnormalized
prepared charged line, and let $\phi$ be its atomic comparison vector.
The complete spectral-projector and preparation bounds give
$\|\chi-\phi\|\leq\epsilon$, $H_0\phi+4N_{\rm ref}\phi=a\phi$,
and $\Pi_0\phi=\phi$, where $\Pi_0$ requires all nonorigin matter cells
to be gauge singlets.  If $w=\|\phi\|^2\in[w_-,w_+]$ and
$p_-=(\sqrt{w_-}-\epsilon)^2>0$, Duhamel's identity gives
\[
 \Pr(\Pi_0\mid\text{line},t)\geq
 1-\min\left\{1,
 \frac{(\epsilon+|t|v\sqrt{w_+})^2}{p_-}\right\}.
\]
The complementary effect includes every screened cloud, as well as
every remote single charge.  For the native $e^c$ removal channel the
origin probability exceeds $0.9999999127$ at preparation and
$0.9995481684$ at $t=10^6$ in source units; the latter interval exceeds
$110{,}036$ atomic source-gap cycles.  These are bounds on the complete
prepared state under its actual Hamiltonian.  They do not apply to every
state in the same energy bin, and the fixed compensator is part of the
construction.

The source's central spin rotation has an exact finite readout on these
same states.  Its native generator is
$S_z=(N_\uparrow-N_\downarrow)/2$, hence
$e^{-2\pi i S_z}=(-1)^N$ on the full source Fock space.  The prepared
neutral cube is even, each successful elementary removal is odd, and
the full Hamiltonian and energy filter preserve that parity.  Controlled
source rotation and a retained coherent pointer therefore distinguish
the minus sign at $2\pi$ from the identity at $4\pi$.  Rotating the odd
compensating reference as well cancels the sign and supplies a distinct
control.  All preparation, pump and energy-bin failures are retained.
This is a finite source and apparatus result; it does not identify a
continuum electron angular momentum or derive an apolar tensor field.
